\section{Full outside protection on a four-vertex path}
\label{graph:section}

This section gives a structured unit-profit graph guarantee. For a finite
resource graph $G$, let
\begin{equation}\label{graph:production}
 F_G(c)=\max\left\{\sum_{e\in E(G)}x_e:
 x_e\ge0,\quad\sum_{e\ni v}x_e\le c_v\ (v\in V(G))\right\}.
\end{equation}
Every edge consumes one unit of each endpoint resource and produces one
unit of output. Owners and their finite coordinatewise endowment boxes
remain fixed; all real reports in those boxes are allowed. Owners may
hold arbitrary mixtures of resources, and every coalition, including its
unchanged members, receives coverage and full outside-production protection.

\begin{theorem}[Protected-leaf path rule]\label{graph:theorem}
For the path $B-A-D-C$, there is an explicit nonnegative payment rule that
satisfies coalition coverage and full OUT for every coalition at every
real report, with
\begin{equation}\label{graph:funding-bound}
 \sum_i y_i(b)\le(1+1/\sqrt2)F_G\left(\sum_i b_i\right).
\end{equation}
The same guarantee holds for complete blow-ups of this path, their finite
disjoint unions, and additional isolated resources, with the same fixed
mixed owners. Thus the total-funding factor is $1+1/\sqrt2$, and the
extra subsidy is at most $F_G(\sum_i b_i)/\sqrt2$.
\end{theorem}

A complete blow-up replaces each of the four vertices by a nonempty
independent group, includes every edge between groups corresponding to a
path edge, and includes no other edges. The theorem is restricted to unit
profits and these graph classes. It asserts neither an all-graph bound
below two nor optimality among payment rules on the path.

\subsection{The selected flow and its payment rule}

Write the path capacities as $c=(A,B,C,D)$, using a vertex label also for
its capacity, and put
\[
 \alpha=\frac1{\sqrt2},\qquad \beta=\alpha^2=\frac12.
\]
These are the constants of this section. For the path,
\begin{equation}\label{graph:path-value}
 F_G(c)=\min\{A+C,B+D,A+D\}.
\end{equation}
Indeed, each expression is a vertex-cover capacity and hence an upper
bound. If $P$ is their minimum, the edge quantities
\[
 x_{AB}=(P-D)_+,\qquad x_{DC}=(P-A)_+,\qquad
 x_{AD}=P-(P-D)_+-(P-A)_+
\]
are feasible and sum to $P$. Here $t_+=\max\{t,0\}$;
nonnegativity of $x_{AD}$ follows from $P\le A+D$, and the other capacity
constraints follow from $P\le A+C$ and $P\le B+D$.

Form a directed bipartite double cover: a left copy of each vertex has
capacity $c_v$, a right copy has capacity $\alpha c_v$, and an edge $uv$
gives arcs $u_Lv_R$ and $v_Lu_R$. Its two components have maximum values
\begin{equation}\label{graph:oriented-values}
 V=\min\{A+C,\alpha(B+D),A+\alpha D\},\qquad
 V'=\min\{B+D,\alpha(A+C),D+\alpha A\}.
\end{equation}
We require a particular choice of left uses, rather than an arbitrary
maximum flow.

\begin{lemma}[All maximum-flow left uses]\label{graph:flow-interval}
In the component $\{A_L,C_L\}\to\{B_R,D_R\}$, the possible left uses
of a maximum flow are exactly
\begin{equation}\label{graph:first-interval}
 w_A+w_C=V,\qquad
 \max\{0,V-A\}\le w_C\le\min\{C,\alpha D,V\}.
\end{equation}
In the other component they are exactly
\begin{equation}\label{graph:second-interval}
 w_B+w_D=V',\qquad
 \max\{0,V'-D\}\le w_B\le\min\{B,\alpha A,V'\}.
\end{equation}
Both intervals are nonempty, including at zero capacities.
\end{lemma}
\begin{proof}
The three bounds defining $V$ are the total left capacity, total right
capacity, and the combined capacity of $A_L,D_R$. They upper-bound every
flow. For any value-$V$ flow, $w_C\le C$, $w_C\le\alpha D$ because
$C_L$ can only send to $D_R$, and $0\le w_C\le V$.
The bound $w_A\le A$ gives the lower endpoint in
\eqref{graph:first-interval}. The interval is nonempty because
$V\le A+C$, $V\le A+\alpha D$, and all capacities are nonnegative.

Conversely, for any $z$ in this interval, take
\[
 f_{A_LB_R}=(V-\alpha D)_+,\qquad
 f_{A_LD_R}=V-z-(V-\alpha D)_+,\qquad f_{C_LD_R}=z.
\]
These quantities are nonnegative: if $V>\alpha D$, use $z\le\alpha D$;
otherwise use $z\le V$. The $B_R$ use is at most $\alpha B$ since
$V\le\alpha(B+D)$, and the $D_R$ use is
$V-(V-\alpha D)_+\le\alpha D$. Their left uses are $V-z\le A$ and
$z\le C$. This realizes every point of the interval and proves that $V$
is the maximum. Reflecting the path, that is, exchanging $A$ with $D$ and
$B$ with $C$, proves the second assertion. Explicitly, with $z=w_B$ its
flows are $(V'-\alpha A)_+$ on $D_LC_R$, $V'-z-(V'-\alpha A)_+$ on
$D_LA_R$, and $z$ on $B_LA_R$.
\end{proof}

Define $\operatorname{clamp}(x;\ell,h)=\min\{h,\max\{\ell,x\}\}$ and
select
\begin{align}
 w_C&=\operatorname{clamp}\bigl(\beta C;
             \max\{0,V-A\},\min\{C,\alpha D,V\}\bigr),
 &w_A&=V-w_C,\label{graph:selector}\\
 w_B&=\operatorname{clamp}\bigl(\beta B;
             \max\{0,V'-D\},\min\{B,\alpha A,V'\}\bigr),
 &w_D&=V'-w_B.\nonumber
\end{align}
The vector $w$ is continuous, nonnegative, homogeneous, and piecewise
linear, with $0\le w_v\le c_v$. The following equivalent formulas expose
its linear cells:
\begin{align}
 w_A&=\min\{A,\alpha B+(\alpha D-\beta C)_+\},
 &w_D&=\min\{D,\alpha C+(\alpha A-\beta B)_+\},\label{graph:explicit-uses}\\
 w_C&=\min\{C,\alpha D,\max\{\beta C,\alpha(B+D)-A\}\},
 &w_B&=\min\{B,\alpha A,\max\{\beta B,\alpha(A+C)-D\}\}.\nonumber
\end{align}
To verify the first orientation algebraically, put
$U=\min\{C,\alpha D\}$, $T=\alpha(B+D)$, $t=\beta C$, and
$m=\min\{U,t\}=\min\{\alpha D,\beta C\}$.
Then $T\ge U$ and $V=\min\{T,A+U\}\ge U$, so the upper endpoint is
$U$. Since $V-A\le U$, the clamp equals
\[
 w_C=\min\{U,\max\{t,V-A\}\}
     =\max\{m,V-A\}
     =\min\{U,\max\{t,T-A\}\}.
\]
Consequently
$w_A=\min\{A,V-m\}=\min\{A,T-m\}$, giving
\eqref{graph:explicit-uses}. Reflection gives the other two formulas.

Set
\begin{equation}\label{graph:prices}
 p_v(c)=\begin{cases}w_v(c)/c_v,&c_v>0,\\0,&c_v=0,\end{cases}
 \qquad y_i(b)=\sum_{v\in\{A,B,C,D\}}b_{iv}p_v\left(\sum_j b_j\right).
\end{equation}
Prices lie in $[0,1]$. The price chosen at an absent coordinate has no
effect on any payment.

The selected left uses are realized by the flows in
Lemma~\ref{graph:flow-interval}. Collapse both directed copies of each
path edge into one edge flow. At vertex $v$, its incident use is at most
$c_v+\alpha c_v$, from the left and right copies respectively. Dividing
by $1+\alpha$ yields a feasible original packing. Therefore
\begin{equation}\label{graph:flow-funding}
 \sum_i y_i(b)=\sum_v w_v(c)=V+V'\le(1+\alpha)F_G(c).
\end{equation}
This argument uses fractional flows and no integrality assumption.

\subsection{The restoration inequalities on every smooth cell}

At a strictly positive aggregate $c$, let $a$ be a coalition's current
aggregate holding, with $0\le a\le c$, and write
$\theta_v=a_v/c_v\in[0,1]$. On a smooth cell define
\begin{equation}\label{graph:restoration-coefficients}
 g_v=p_v+\sum_j a_j\partial_v p_j.
\end{equation}
This differentiates the coalition's total payment when its holding and
the aggregate capacity are restored together, while outsiders stay fixed.
We prove
\begin{equation}\label{graph:cover-certificate}
 g_v\ge0\quad\text{for every vertex},\qquad
 g_A+g_B\ge1,\quad g_A+g_D\ge1,\quad g_D+g_C\ge1.
\end{equation}
Ties and zero capacities will be handled by direct integration below.

First record the single-edge calculation used in the first cells. For
edge capacities $X\ge Y>0$, put $t=Y/X$ and take
$w_X=\alpha Y$, $w_Y=\min\{Y,\alpha X\}$.
When $t\le\alpha$, the restoration coefficients are
\[
 g_X=\alpha t(1-\theta_X),\qquad g_Y=1+\alpha\theta_X.
\]
They are nonnegative and their sum is
$1+\alpha t+\alpha(1-t)\theta_X\ge1$.
When $t\ge\alpha$, they are
\[
 g_X=\alpha t(1-\theta_X)+\alpha\theta_Y,\qquad
 g_Y=(\alpha/t)(1-\theta_Y)+\alpha\theta_X.
\]
They are nonnegative. Minimizing their sum over the holding fractions gives
$\alpha(1+t)\ge\alpha(1+\alpha)>1$: the coefficient of $\theta_X$ in
the sum is nonnegative, and that of $\theta_Y$ is nonpositive.
Swapping $X,Y$ handles the other capacity ordering.

For the path put
\[
 u=A-\alpha B,\qquad v=D-\alpha C.
\]
The following cases exhaust the smooth cells; the displayed weak
inequalities describe their closures.

\paragraph{A nonpositive difference.}
If $u\le0$, formula \eqref{graph:explicit-uses} reduces to
\[
 w_A=A,\quad w_B=\alpha A,\quad
 w_C=\min\{C,\alpha D\},\quad w_D=\min\{D,\alpha C\}.
\]
For the first pair,
\[
 g_A=1+\alpha\theta_B\ge1,\qquad
 g_B=(\alpha A/B)(1-\theta_B)\ge0.
\]
The other pair satisfies the single-edge calculation, so $g_C,g_D\ge0$
and $g_C+g_D\ge1$. The cross edge is covered by $g_A\ge1$.
If $v\le0$, reflection gives the same conclusions with $g_D\ge1$.

For the remaining cells $u,v>0$, and the common formulas are
\begin{align*}
 w_A&=\alpha B+\min\{u,\alpha v\},&
 w_D&=\alpha C+\min\{v,\alpha u\},\\
 w_B&=\min\{B,\beta B+(\alpha u-v)_+\},&
 w_C&=\min\{C,\beta C+(\alpha v-u)_+\}.
\end{align*}
Since $\alpha<1$, at most one of $\alpha u-v$ and $\alpha v-u$ is
positive.

\paragraph{The middle cone: $\alpha v\le u$ and $\alpha u\le v$.}
Here
\[
 w_B=\beta B,\quad w_C=\beta C,\quad
 w_A=\alpha(B+D)-\beta C,\quad w_D=\alpha(A+C)-\beta B.
\]
Moreover $p_A,p_D\ge\beta$. For example,
$\alpha u\le v$ gives
$w_A=\alpha B+\alpha v\ge\beta(\alpha B+u)=\beta A$;
the other inequality follows by reflection. Direct differentiation yields
\begin{align*}
 g_A&=p_A(1-\theta_A)+\alpha\theta_D,&
 g_D&=p_D(1-\theta_D)+\alpha\theta_A,\\
 g_B&=\beta+\alpha\theta_A-\beta\theta_D,&
 g_C&=\beta+\alpha\theta_D-\beta\theta_A.
\end{align*}
All are nonnegative. Their edge-sum minima over $\theta\in[0,1]^4$ are
\[
 \min(g_A+g_B)=\beta+\min\{p_A,\alpha\},\qquad
 \min(g_D+g_C)=\beta+\min\{p_D,\alpha\},
\]
and
\[
 \min(g_A+g_D)=\min\{p_A,\alpha\}+\min\{p_D,\alpha\}.
\]
Every displayed value is at least $2\beta=1$.

\paragraph{An outer cone: $u\le\alpha v$.}
The opposite outer cone is its reflection. In this cone
\[
 w_A=A,\quad w_B=\beta B,\quad
 w_D=\alpha(A+C)-\beta B,\quad
 w_C=\min\{C,\alpha(B+D)-A\}.
\]
If $w_C<C$, then $p_C\ge\beta$ because $\alpha v-u\ge0$.
The unsaturated condition gives
$\alpha v-u\le(1-\beta)C$. Multiplying by $\alpha$ and using
$\beta=\alpha^2$ gives
\[
 \beta v\le\alpha u+\alpha(1-\beta)C,
 \qquad \beta D\le\alpha u+\alpha C=w_D.
\]
Thus $p_D\ge\beta$ as well. The coefficients are
\begin{align*}
 g_A&=1-\theta_C+\alpha\theta_D,&
 g_B&=\beta+\alpha\theta_C-\beta\theta_D,\\
 g_C&=p_C(1-\theta_C)+\alpha\theta_D,&
 g_D&=p_D(1-\theta_D)+\alpha\theta_C.
\end{align*}
They are nonnegative. Direct affine minimization gives
\begin{align*}
 \min(g_A+g_B)&=\alpha+\beta\ge1,\\
 \min(g_C+g_D)&=\min\{p_C,\alpha\}+\min\{p_D,\alpha\}\ge2\beta=1,\\
 \min(g_A+g_D)&=\alpha+\min\{p_D,\alpha\}\ge\alpha+\beta\ge1.
\end{align*}
For example the first sum is
$1+\beta+(\alpha-1)\theta_C+(\alpha-\beta)\theta_D$, whose minimum
is $\alpha+\beta$.

If instead $w_C=C$ on a smooth saturated cell, the coefficients are
\[
 g_A=1+\alpha\theta_D,\qquad g_C=1+\alpha\theta_D,
 \qquad g_B=\beta(1-\theta_D),\qquad
 g_D=p_D(1-\theta_D).
\]
They are nonnegative; $g_A\ge1$ covers $AB$ and $AD$, and $g_C\ge1$
covers $CD$. This completes \eqref{graph:cover-certificate} on every
smooth cell for every holding fraction, without restricting mixed ownership.

\subsection{Finite restorations, ties, and zero reports}

For outsider aggregate $o\ge0$ and coalition aggregate $a\ge0$, define
\[
 H(o,a)=\sum_v a_v p_v(o+a).
\]
This function is jointly continuous on its entire nonnegative domain.
At a coordinate whose aggregate tends to zero, its contribution lies
between zero and $a_v$ and hence tends to zero; at a positive aggregate
coordinate, continuity follows from that of $w$. Also $H(o,0)=0$.

Fix $0\le a'\le a$ and put $d=a-a'$. The finitely many polyhedral cells
of $w$ have boundaries contained in finitely many proper hyperplanes.
Choose a strictly positive vector $z$ having nonzero inner product with
each of their nonzero normals. Such a vector exists because a finite union
of proper hyperplanes cannot cover the positive orthant. Choose
$\varepsilon\downarrow0$ omitting the finitely many positive values for
which the segment
\[
 c(t)=o+\varepsilon z+a'+td,\qquad 0\le t\le1,
\]
is contained in any boundary hyperplane. Indeed, containment requires the
hyperplane to be parallel to $d$ and imposes one linear equation in
$\varepsilon$ with a nonzero coefficient. All other boundary intersections
occur at at most one parameter value per hyperplane.

The segment is strictly positive. Its coalition payment
$H(o+\varepsilon z,a'+td)$ is continuous and piecewise continuously
differentiable, with only finitely many breakpoints. Away from them its
derivative is $g\cdot d$, and its holding fractions are
$(a'_v+td_v)/c_v(t)\in[0,1]$. The verified inequalities therefore apply.
For any edge packing $x$ feasible at $d$, they imply
\[
 \sum_e x_e\le\sum_{e=uv}(g_u+g_v)x_e
 =\sum_v g_v\sum_{e\ni v}x_e\le\sum_v g_vd_v.
\]
Taking the maximum gives $g\cdot d\ge F_G(d)$. Integrating on the
finitely many pieces, using continuity at their endpoints, yields
\[
 H(o+\varepsilon z,a)-H(o+\varepsilon z,a')\ge F_G(d).
\]
Joint continuity as $\varepsilon\downarrow0$ proves
\begin{equation}\label{graph:finite-out}
 H(o,a)-H(o,a')\ge F_G(a-a')
 \qquad(o\ge0,\ 0\le a'\le a).
\end{equation}
The perturbation is an analytic approximation to the globally defined
formula $H$; no perturbed owner profile is required to fit the fixed boxes,
and no owner is added or reassigned in the asserted mechanism.

For an actual coalition $T$, take $o=\sum_{i\notin T}b_i$,
$a=\sum_{i\in T}b_i$, and $a'=\sum_{i\in T}b'_i$ for its binding
underreports. Its aggregate payment is exactly $H(o,a)$, so
\eqref{graph:finite-out} is full OUT for the original owners, including
passive members. It includes segments on ties, zero capacities, the
all-zero profile, and the grand coalition. Taking $a'=0$ proves coverage.
Together with \eqref{graph:flow-funding}, this proves
Theorem~\ref{graph:theorem} for the path.

\subsection{Complete blow-ups and disjoint unions}

Let $X\in\{A,B,C,D\}$ now denote one of four independent vertex groups,
and aggregate capacities as $c_X=\sum_{v\in X}c_v$. The production value
of a complete blow-up equals the path value at these four aggregates.
Aggregating any original packing proves one inequality. To prove the
other, let $f_{XY}$ be a feasible aggregate path packing. For $c_X>0$ assign
the incidence of this group-pair flow to $v\in X$ by
\[
 r_{v,XY}=\frac{c_v f_{XY}}{c_X}.
\]
The sum over incident group pairs is at most $c_v$, and
$\sum_{v\in X}r_{v,XY}=f_{XY}$. If $f_{XY}>0$, both group capacities
are positive and we may put
\[
 x_{vw}=\frac{r_{v,XY}r_{w,XY}}{f_{XY}}
 \qquad(v\in X,\ w\in Y)
\]
on every edge of that complete bipartite pair. Its two marginals are the
assigned incidences, and its total is $f_{XY}$. Zero-flow pairs receive
zero flow. Zero-capacity groups have no positive incident aggregate flow,
so no division by zero is used. This constructs an original feasible
packing of exactly the aggregate objective.

Apply the four path prices to every resource in its corresponding group.
Each actual owner's payment equals the path-rule payment on its own
aggregated report. Every actual withholding maps to a coordinatewise
withholding of that same owner's aggregate, and the equality of production
values holds for all coalition and withheld bundles. Thus funding,
coverage, and full OUT transfer exactly, retaining the original mixed
owners and their fixed endowment boxes.

For a finite disjoint union, sum the component payments. Both production
and withheld production are additive across components, so summing their
OUT, coverage, and funding inequalities preserves the constant $1+\alpha$.
An owner may hold resources in several components. Isolated resources
contribute no production and may receive zero payment. This completes
Theorem~\ref{graph:theorem}.

\begin{remark}[Scope of the funding number]\label{graph:scope}
At $(A,B,C,D)=(4,1,1,1)$, the path value is two and
$(V,V')=(2\alpha,2)$, so this particular rule attains the ratio
$1+\alpha$. This is not a lower bound for arbitrary mechanisms on the
path. The proof does not extend the guarantee to arbitrary edge profits
or to incomplete connections between the four groups, and makes no
claim of historical novelty.
\end{remark}
